\section{Introdução}

A incerteza é um elemento inerente a muitos fenômenos estudados na Ciência da
Computação, desde o tempo de resposta de um servidor até o resultado de um
classificador automático. A teoria da probabilidade oferece a linguagem formal
para descrever esses fenômenos e raciocinar sobre eles de maneira quantitativa,
tendo se tornado indispensável na Computação \cite{harcholbalter2024probability}.
Sua base axiomática moderna foi estabelecida por Andrey Kolmogorov em 1933
\cite{kolmogorov1933grundbegriffe}. Na engenharia, a mesma teoria fundamenta a
metodologia estatística empregada na resolução de problemas reais
\cite{montgomery2018applied}.

Boa parte da Computação atual gira em torno de dados, que são inseparáveis da
probabilidade \cite{chan2021probability}. A estatística reúne os métodos de coleta,
análise e interpretação de dados que permitem extrair informação e apoiar decisões
\cite{sahu2024probability}, de modo que estatística e computação se consolidam como
recursos fundamentais na pesquisa e na indústria \cite{ferreira2020estatistica}.
Conceitos como esperança, variância, independência e o Teorema de Bayes aparecem de
forma recorrente em problemas de classificação, diagnóstico e avaliação de
desempenho.

Nem sempre as probabilidades de interesse podem ser obtidas por via puramente
analítica, e a simulação computacional surge como alternativa para estimá-las pela
repetição de um experimento aleatório. Essa abordagem é amplamente usada na
avaliação do desempenho de sistemas computacionais
\cite{harcholbalter2024probability} e tem respaldo na Lei dos Grandes Números, que
garante a convergência das frequências observadas para as probabilidades teóricas
\cite{ross2010probabilidade}. As simulações deste trabalho foram conduzidas em R,
uma linguagem e ambiente para computação estatística e gráfica
\cite{ferreira2020estatistica}, bastante usado para ilustrar conceitos de
probabilidade como o próprio Teorema Central do Limite \cite{sahu2024probability}.
Em todos os experimentos foi fixada a semente do gerador de números aleatórios, o
que assegura a reprodutibilidade dos resultados.

O objetivo deste relatório é revisitar conceitos centrais da probabilidade
combinando o tratamento teórico com experimentos computacionais em R, de modo que
cada propriedade enunciada seja também verificada de forma empírica. O texto está
organizado em torno desses dois eixos, o teórico e o computacional. A primeira
parte trata da simulação de experimentos aleatórios simples, como o lançamento de
uma moeda, a retirada de bolas de uma urna e a soma de dois dados, e introduz a
frequência relativa como estimador de probabilidade. Em seguida são estudadas as
distribuições discretas e as distribuições contínuas, com destaque para o Teorema
Central do Limite, que descreve a aproximação da média amostral à distribuição
normal. A esperança e a variância são então analisadas em conjunto com a Lei dos
Grandes Números, que formaliza a estabilidade das médias amostrais. A parte
seguinte aborda a probabilidade condicional, a independência de eventos e o Teorema
de Bayes, aplicados a problemas como filtragem de spam, diagnóstico médico e
autenticação biométrica. Por fim, uma atividade integradora final reúne vários
estudos que aplicam e combinam os conceitos anteriores em problemas de inspiração
computacional, alguns deles apoiados em modelos clássicos da probabilidade aplicada
\cite{ross2024models}.
